% !TeX root = ../教材题库_全部条目_转写工作稿.tex
% Source: XJ-B1-C02-TH-P040-Y0846
% 第2章 / 2.1 相等关系与不等关系 / 教材 p.40、41 / PDF 第47、48页
% OCR未校对：公式与文字均可能有误；未完成数学TeX转写；禁止直接发给学生
% 原题及相邻原解截图（核对用）：../题目原图/XJ-B1-C02-TH-P040-Y0846-01.webp, ../题目原图/XJ-B1-C02-TH-P040-Y0846-02.webp
% 本文件只在转写完成后作为题干使用；OCR 条目尚不是可靠题干。
\begin{bankitem}{p.40 思考（含紧邻说明）}{第2章 / 2.1 相等关系与不等关系 / 教材 p.40、41 / PDF 第47、48页}{XJ-B1-C02-TH-P040-Y0846}{OCR未校对：公式与文字均可能有误；未完成数学TeX转写；禁止直接发给学生}
\begin{ocrdraft}
[教材p.40 / PDF第47页]
\par
？思考（D把12写成两个正数的乘积，当这两个正数取什么值时，它们的和
\par
最小？
\par
（2）把25写成两个正数的和，当这两个正数取什么值时，它们的积最大？
\par
（3）由（1）（2）你能概括出什么结论？
\par
（1)设两个正数为则>0，y>0，且一12．
\par
由
\par
2
\par
可得
\par
十y冫2《0，一212一4，/S，
\par
工+，冫4而，就说明
\par
当且仅当一3，时等号成立，此时一3，一2．
\par
工+，始终大于或等于
\par
所以，把12写成两个2的乘积时，它们的和最小，最
\par
4，取等号时，两数之
\par
和最小．
\par
小和为4．
\par
（2）设两个正数为产则>0，3>0，且十3，一25．
\par
由
\par
2
\par
可得
\par
625
\par
4
\par
[教材p.41 / PDF第48页]
\par
当且仅当工一，时等号成立，此时=3-
\par
625
\par
所以，把25写成两个的和时，它们的积最大，最大积为．一一
\par
（3）由（1）（2）可概括出下述结论：
\par
已知3，都为正数，则
\par
（1）如果积是定值，那么当且仅当x=y时，和x十3，有最小值2，/
\par
0
\par
（2）如果和x十3，是定值5，那么当且仅当一时，积《y有最大值\_
\par
数学上已经证明上述结论成立．
\par
下面结合具体实际问题，感受上述结论在解决最大（小）值问题中的作用．
\end{ocrdraft}
\end{bankitem}
